
% ============================================================================
\clearpage
\section*{v4.10.0-r61 Source-Only: Corrections to r56.4 (Normalization of \texorpdfstring{\(\rho_\Lambda\)}{rho Lambda}) and r07.4 (Square Root in \texorpdfstring{\(\beta_{\rm sym}\)}{beta sym}) (k320)}
\addcontentsline{toc}{section}{v4.10.0-r61 Source-Only: Corrections to r56.4 (Normalization of \texorpdfstring{\(\rho_\Lambda\)}{rho Lambda}) and r07.4 (Square Root in \texorpdfstring{\(\beta_{\rm sym}\)}{beta sym}) (k320)}
\label{sec:v4100-r61-corrections}
% ============================================================================

\paragraph{Purpose.}
While the English article ``One Object, Three Frames'' (\texttt{paper/\allowbreak physics-standalone/}) was checked against the Companion on
25 November 2026, two formula slips were found in earlier layers.  On the author's request this layer corrects them.  Earlier layers are not
rewritten; the corrections below replace the two displayed formulas wherever r56.4 and r07.4 are used.  Both are corrections of
presentation: no number, no test result and no status in the record changes.  This layer is an exploration and \emph{not blind}; no hypothesis
number is assigned.

\noindent\textbf{Source-only status.}
This v4.10.0-r61 layer is source-only.  It does not build a PDF, does not assign a new DOI,
and does not authorize a public release.  The v5.4.0 Version DOI
\href{https://doi.org/10.5281/zenodo.24500009}{10.5281/zenodo.24500009} and the Concept DOI
\href{https://doi.org/10.5281/zenodo.19210728}{10.5281/zenodo.19210728} are unchanged by it.

\paragraph{Theoremlet v4.10.0-r61.1 (correction to r56.4: normalization of the vacuum density).}
Exploration, checked with Wolfram Language (\texttt{k320}).  In flat FLRW with \(H^2=\kappa\rho/3\), \(\kappa=8\pi G\), a cosmological constant
contributes \(\rho_\Lambda=\Lambda/\kappa\), so that \(H^2=\kappa\rho_M/3+\Lambda/3\); this is also the form used in r56.2 (\texttt{k242}, part A).
r56.4 wrote \(\rho_\Lambda=3\Lambda/\kappa\).  With that normalization the relation to the Hubble rate would be \(\Lambda=(1-\beta^2)H_*^2\), not the
stated \(\Lambda=3(1-\beta^2)H_*^2\); the two displayed formulas of r56.4 do not hold together.  With \(\rho_\Lambda=\Lambda/\kappa\) the condition
\(\Omega_M(x)=\beta^2\) has for every \(\Lambda>0\) exactly one positive solution,
\[
x_*^3=\frac{\kappa\rho_{M0}(1-\beta^2)}{\beta^2\Lambda},\qquad \Lambda=3(1-\beta^2)H_*^2=2.0944603\,H_*^2 .
\]
Only the expression for \(x_*^3\) changes, by the factor 3 in its denominator.  The second formula, its number \(2.0944603\), the uniqueness of the
epoch and the exchange of \(\Lambda\) and epoch under the readout H96 stand as stated.  Corollary r56.5 and the later uses of r56.4 (r57.1, r59.5,
\texttt{k251}) rely only on these and are unchanged.  \texttt{k242}, part C, used the normalization of r56.4 for \(x_*\); its printed relation
\(\Lambda=3(1-\beta^2)H_*^2\) was computed with \(\Omega_\Lambda=1-\beta^2\) directly and is correct.

\paragraph{Theoremlet v4.10.0-r61.2 (correction to r07.4: square root in \(\beta_{\rm sym}\)).}
Exploration, checked with Wolfram Language (\texttt{k320}).  r07.4(c) wrote \(\beta_{\rm sym}=(\gamma_L-1)/(\gamma_L+1)\).  The symmetric-frame
offset is
\[
\beta_{\rm sym}=\sqrt{\frac{\gamma_L-1}{\gamma_L+1}}=0.5494056,\qquad \beta_{\rm sym}^2=\frac{\gamma_L-1}{\gamma_L+1}=0.3018466
\]
(\texttt{thm:symmetric-frame}).  Without the root the expression is \(\beta_{\rm sym}^2\), and ``\(\Omega_M=\beta_{\rm sym}^2\)'' read literally from
r07.4 would give \(0.0911\) instead of the value \(0.30185\) of H96.  No number in the record was computed from the formula as written in r07.4;
all values of \(\beta_{\rm sym}\), \(\beta_{\rm sym}^2\) and H96 use the root.  The statement of r07.4(c) is unaffected: every function of
\(\gamma_L=\kappa_{\rm DL/M}/\kappa_{\rm ABCK}\), with or without the root, is independent of the area-law prefactor \(N\) (\texttt{k320}, part E).

\subsection*{v4.10.0-r61 what depends on the corrections}
\label{subsec:v4100-r61-impact}

\begin{center}
\small
\begin{tabular}{@{}>{\raggedright\arraybackslash}p{0.34\textwidth}>{\raggedright\arraybackslash}p{0.58\textwidth}@{}}
\toprule
place & effect\\
\midrule
r56.4, formula for \(x_*^3\) & replaced by r61.1 (factor 3)\\
r56.4, \(\Lambda=3(1-\beta^2)H_*^2\); r56.5; r57.1; r59.5 & unchanged\\
H96, H124, all values of \(\beta_{\rm sym}\) & unchanged\\
r07.4(c), formula for \(\beta_{\rm sym}\) & replaced by r61.2 (square root); conclusion unchanged\\
\texttt{k242} part C & normalization of \(x_*\) as in r56.4; its \(\Lambda/H_*^2\) value correct\\
physics paper v5.4.0, preprint 0.3.0 & neither formula quoted; no change needed\\
\bottomrule
\end{tabular}
\end{center}

\subsection*{v4.10.0-r61 frames}
\label{subsec:v4100-r61-frames}

\begin{center}
\small
\begin{tabular}{@{}lll@{}}
\toprule
item & frame & comparison\\
\midrule
r61.1 (cosmological readout) & M (matter-frame observer, as r56.4) & none across frames\\
r61.2 (\(\beta_{\rm sym}\), counting ratio) & M and T through \(\gamma_L\), as r07.4 & unchanged\\
P & not used & ---\\
\bottomrule
\end{tabular}
\end{center}

\subsection*{v4.10.0-r61 notes}
\label{subsec:v4100-r61-notes}

{\raggedright
\texttt{notes/\allowbreak research/\allowbreak v4\_\allowbreak 10\_\allowbreak 0\_\allowbreak r61\_\allowbreak korrekturen/\allowbreak llm/\allowbreak k320\_\allowbreak korrektur\_\allowbreak r56\_\allowbreak r07.wl}
with its actual output \texttt{k320\_\allowbreak out.txt}.
\par}

\subsection*{v4.10.0-r61 ledger and guard update}
\label{subsec:v4100-r61-ledger-guard}

\[
R_{\rm chk}^{v4100\text{-}r61}=(1,1,0)
\]
for (corrections recorded; computations checked in \texttt{k320}; no new physical statement).

The layer spends no guard:
\[
R_{\rm glob}^{(v4100\text{-}r61,+)}=(0,1,1,1,1),
\qquad
\Gamma_{\rm guard}^{v4100\text{-}r61}=(1,1,1,1,1),
\qquad
\beta_4=1/256\ \text{unspent}.
\]
The layer does not imply exact finite \(C5\) zero closure, exact all-readout terminality, \(C6/T2\), a global photon-clock/RH theorem,
or global all-readout closure.  OP1--OP5, C5 and C6/T2 remain open.
